\section{Lerch deformation and finite polylogarithmic closure}
\label{ag:sec:polylogs}

For $|z|<1$, replace the denominator in \eqref{ag:eq:Hmellin} by
$1-ze^{-t}$ and define
\begin{equation}\label{ag:eq:Lerch-transform}
 \HH_r(s;a,b;z)=\frac1{\Gamma(s)}\int_0^\infty
             \frac{t^{s-1}e^{-at}g_r(bt)}{1-ze^{-t}}\dd t,
                     \qquad\Re s>0.
\end{equation}
Let $\Phi(z,s,x)=\sum_{n\ge0}z^n(n+x)^{-s}$ in the right half-plane.
The Lerch series, its Mellin formula, and its relation to polylogarithms
are classical \cite{DLMFlerch}. Termwise spatial differentiation gives
$\partial_x\Phi(z,s,x)=-s\Phi(z,s+1,x)$.

\begin{theorem}[Lerch and polylogarithmic endpoints]\label{ag:thm:Lerch}
For $|z|<1$,
\begin{align}
 \HH_r(s;a,b;z)&=\frac{c_r}{2}\int_{-1}^1
                       (1-v^2)^r\Phi(z,s,a+bv)\dd v,
                       \label{ag:eq:Lerch-compact}\\
 &=\sum_{j=r}^{2r}w_{rj}(b)
 \frac{\Phi(z,s-j-1,a+b)-(-1)^j\Phi(z,s-j-1,a-b)}{(1-s)_{j+1}}.
                       \label{ag:eq:Lerch-endpoint}
\end{align}
The function is entire in $s$. Hence every apparent integer pole on
the right is removable, including $s=1$.
If $a\pm b$ are positive integers, every endpoint is a finite
polylogarithmic expression, using
\begin{equation}\label{ag:eq:integer-Lerch}
 \Phi(z,s,N)=z^{-N}\left(\Li_s(z)-\sum_{n=1}^{N-1}\frac{z^n}{n^s}\right)
 \quad(N\ge1,
 z\ne0).
\end{equation}
All values at $z=0$ are supplied by analytic limits.
\end{theorem}
\begin{proof}
The beta representation and absolute convergence prove
\eqref{ag:eq:Lerch-compact}. The endpoint argument in
Theorem~\ref{ag:thm:hurwitz} applies unchanged. Exponential decay in
$n$ makes the Lerch series locally uniformly convergent for every
$s$, including all spectral derivatives. This proves entire continuation
and removability of the displayed poles. Shifting the summation index
proves \eqref{ag:eq:integer-Lerch}.
\end{proof}

The limit $z\to1$ is initially justified for $\Re s>1$, where it
recovers the Hurwitz transform. Continuing in $s$ before taking that
limit need not be interchangeable at $s=1$: the Hurwitz pole must be
kept. At a removable integer in \eqref{ag:eq:Lerch-endpoint}, the value
may involve \emph{order derivatives} of polylogarithms; a formula
containing a vanishing denominator should not be read as an ordinary
substitution rule.

\subsection{Rational spatial endpoints}
For $0<z<1$, let $w=z^{1/q}$, $\omega=e^{2\pi\ii/q}$, and take
positive integers $A,q$. The roots-of-unity filter in the absolutely
convergent polylogarithm series gives
\begin{align}
 \Phi(z,s,A/q)
 ={}&q^{s-1}w^{-A}\sum_{\ell=0}^{q-1}\omega^{-\ell A}
                                      \Li_s(w\omega^\ell)\notag\\
 &-q^sw^{-A}
       \sum_{\substack{1\le n<A\\n\equiv A\;({\rm mod}\ q)}}
                                    \frac{w^n}{n^s}.
 \label{ag:eq:rational-Lerch}
\end{align}
Indeed the finite Fourier sum retains exactly the positive integers
congruent to $A$ modulo $q$, and the finite prefix removes those below
$A$. The formula holds for every complex $s$ by normal convergence.
Its order derivatives give finite log-weighted polylogarithm reductions
for all spectral derivatives of the rational-endpoint transform.

\subsection{An explicit expression with only three polylogarithms}
Put
\[
 M(t)=(t^2+3)\sinh t-3t\cosh t,
\]
and, for $z\ne0$,
\[
 J_\sigma^\pm(z)=(z^{-3}\pm z^{-1})\Li_\sigma(z)
                                  -z^{-2}-2^{-\sigma}z^{-1}.
\]
For $r=2,a=2,b=1$, \eqref{ag:eq:Lerch-endpoint} reads
\begin{equation}\label{ag:eq:r2-Lerch}
 \HH_2(s;2,1;z)=\frac{J_{s-3}^-(z)}{96(1-s)_3}
       -\frac{J_{s-4}^+(z)}{32(1-s)_4}
       +\frac{J_{s-5}^-(z)}{32(1-s)_5}.
\end{equation}
At $s=6$ there is no exceptional denominator. Since
$g_2(t)=M(t)/(48t^5)$, this gives the ordinary convergent identity
\begin{align}
 \int_0^\infty\frac{e^{-2t}M(t)}{1-ze^{-t}}\dd t
 ={}&\frac4{z^2}+\frac5{4z}
 -\frac{1-z^2}{z^3}\left(\Li_3(z)+\frac32\Li_1(z)\right)\notag\\
 &-\frac{3(1+z^2)}{2z^3}\Li_2(z),\qquad 0<z<1.
 \label{ag:eq:three-polylogs}
\end{align}
The finite limit at $z=1$ is $21/4-\pi^2/2$. The individual negative
powers of $z$ cancel at $z=0$. These endpoint cancellations are most
transparent in the integral or the Lerch series, rather than in the
expanded polylogarithmic expression.
